% language=us

\environment context-2025-style

% \setuptyping
%   [style=\switchtobodyfont[7pt,tt]]

\definelayer
  [extratext]
  [preset=rightbottom,
   height=\textheight,
   width=\textwidth]

\setupbackgrounds
  [text]
  [background=extratext]

% \definehighlight[notabene][style=bf]
\definehighlight[notabene][style=\glyphweight100\underbar]

\setuptyping
  [align=hangright,
   keeptogether=paragraph,
   numbering=]

\startdocument
  [title={overload protection},
   banner={what is it and why is it needed},
   location={context\enspace{\bf 2025}\enspace meeting}]

\starttitle[title=What is it]

\startitemize
\startitem
In principle a user can any moment change the way \TEX works.
\stopitem
\startitem
This can be intentional or by accident.
\stopitem
\startitem
In doing so side effects can get unnoticed.
\stopitem
\startitem
Candidates for change are \TEX\ macros and registers but also \LUA\
and \METAPOST\ code.
\stopitem
\startitem
So how can we prevent harm being done.
\stopitem
\stopitemize

\stoptitle

\starttitle[title=How does it work: macros]

Users can define macros:

\starttyping
\def\MyMacro{do something}
\stoptyping

But what if that one is also a system one?

\starttyping
\permanent\def\MyMacro{do something}
\stoptyping

This tags a macro as one that can't be changed.

\starttyping
\overloadmode=4 % warning
\overloadmode=5 % error
\stoptyping

\stoptitle

\starttitle[title=How does it work: variables]

Users can define variables:

\starttyping
\integerdef\MyInteger 123
\stoptyping

These definitions can also be protected:

\starttyping
\permanent\integerdef\MyInteger 0
\stoptyping

But values can still be changed unless we say:

\starttyping
\immutable\MyInteger 123
\stoptyping

From then on they can't be changed. The same applies to for instance \type
{\countdef}.

The \type {\newinteger} and \type {\newcount} definition macros already
make the target permanent.

\stoptitle

\starttitle[title=Impact on engine]

\startitemize
\startitem
    We need to relate states to all entries in the so called table of
    equivalents.
\stopitem
\startitem
    Quite some properties are set by prefixes; we can group them.
\stopitem
\startitem
    Obsolete: \typ {\long}, \typ {\outer}.
\stopitem
\startitem
    Backend:  \typ {\deferred}, \typ {\immediate}.
\stopitem
\startitem
    Scope: \typ {\global}.
\stopitem
\startitem
    Expansion: \typ {\protected}, \typ {\semiprotected}.
\stopitem
\startitem
    Kind of special: \typ {\noaligned}, \typ {\untraced}, \typ {\constrained},
    \typ {\retained}, \typ {\tolerant}, \typ{\constant}.
\stopitem
\startitem
    Protection: \typ {\aliased}, (\typ {\always}), (\typ {\enforced}), \typ
    {\frozen}, \typ {\immutable}, \typ {\inherited}, \typ {\instance}, \typ
    {\mutable}, \typ {\overloaded}, \typ {\permanent}.
\stopitem

\stopitemize

\stoptitle

\starttitle[title=Candidates for protection]

\startcolumns \small \setupinterlinespace
\starttabulate[|||]
\BC arithmic              \NC \type {\advance}         \NC \NR
\BC internal integer      \NC \type {\hangafter}       \NC \NR
\BC register integer      \NC \type {\countdef}        \NC \NR
\BC internal attribute    \NC \type {(none)}           \NC \NR
\BC register attribute    \NC \type {\attributedef}    \NC \NR
\BC internal posit        \NC \type {(none)}           \NC \NR
\BC register posit        \NC \type {\floatdef}        \NC \NR
\BC internal dimension    \NC \type {\hangindent}      \NC \NR
\BC register dimension    \NC \type {\dimendef}        \NC \NR
\BC font property         \NC \type {\efcode}          \NC \NR
\BC internal glue         \NC \type {\baselineskip}    \NC \NR
\BC register glue         \NC \type {\skipdef}         \NC \NR
\BC internal muglue       \NC \type {\medmuskip}       \NC \NR
\BC register muglue       \NC \type {\muskipdef}       \NC \NR
\BC internal toks         \NC \type {\everypar}        \NC \NR
\BC register toks         \NC \type {\toksdef}         \NC \NR
\BC define char code      \NC \type {\catcode}         \NC \NR
\BC def                   \NC \type {\def}             \NC \NR
\BC define family         \NC \type {\scriptfont}      \NC \NR
\BC define font           \NC \type {\font}            \NC \NR
\BC hyphenation           \NC \type {\hjcode}          \NC \NR
\BC let                   \NC \type {\let}             \NC \NR
\BC prefix                \NC \type {\global}          \NC \NR
\BC register              \NC \type {\count}           \NC \NR
\BC auxiliary             \NC \type {\prevgraf}        \NC \NR
\BC set box               \NC \type {\setbox}          \NC \NR
\BC box property          \NC \type {\wd}              \NC \NR
\BC set font              \NC \type {\nullfont}        \NC \NR
\BC interaction           \NC \type {\batchmode}       \NC \NR
\BC math parameter        \NC \type {\setmathatomrule} \NC \NR
\BC page property         \NC \type {\pagegoal}        \NC \NR
\BC specification         \NC \type {\parshape}        \NC \NR
\BC shorthand def         \NC \type {\chardef}         \NC \NR
\BC association           \NC \type {\associateunit}   \NC \NR
\BC lua value             \NC \type {(in lua)}         \NC \NR
\BC integer               \NC \type {\integerdef}      \NC \NR
\BC posit                 \NC \type {\positdef}        \NC \NR
\BC dimension             \NC \type {\dimensiondef}    \NC \NR
\BC gluespec              \NC \type {\gluespecdef}     \NC \NR
\BC mugluespec            \NC \type {\mugluespecdef}   \NC \NR
\BC index                 \NC \type {\parameterdef}    \NC \NR
\BC combine toks          \NC \type {\toksapp}         \NC \NR
\stoptabulate
\stopcolumns

\blank[2*big]

We need to deal with protection at the \TEX\ end as well as at the \LUA\
end.

In \LUA\ we can to some extend discourage overloading of functionality and
local binding is also an option.

In \METAPOST\ we also have some mechanism but for now we don't error there
(conceptually different grouping model).

% \def\foo#1{\parameterdef\FirstFoo\plusone \FirstFoo}

\stoptitle

\starttitle[title=Virus protection]

Can \TEX\ ever be really secured?

\startitemize
\startitem
    Because \TEX\ reads and writes files, much depends on trust.
\stopitem
\startitem
    It can never be made complete safe but to some extend we can at least try and
    give the user a good feeling.
\stopitem
\startitem
    Can we be sure that  \TEX ies are the \quote {good guys}?
\stopitem
\startitem
    In \CONTEXT\ we can sandbox, that is, all tricky mechanisms have intercepts and
    can be controlled from the configuration. No one ever does that.
\stopitem
\startitem
    The community can impose some restrictions and/or guard integrity of
    macro packages. As it contradicts open source there is no such policy.
\stopitem
\startitem
    Can we intercept injection of potential harmful code? Because one can
    adapt the source early on there is little we can do. So it depends.
\stopitem
\startitem
    To some extend overload protection can help but bad (also \TEX) actors
    can always get around it. Unless we take more drastic measures.
\stopitem
\startitem
    In that respect it is like virus scanners (operating system, browsers,
    apps on phones, etc): the best you can do is observe and adapt.
\stopitem
\stopitemize

\stoptitle

\stopdocument

A \CONTEXT\ meeting is a good moment to consider and discuss how to move on. What
features are needed, which could be improved, and/or should we drop something
that is never used at all? What additional infra structure is needed (helpers,
scripts) and to what extend are users really programming in \TEX.

The \LUAMETATEX\ engine brings quite some new features. Think of more capable
macro arguments, more conditionals, native loops, expressions etc. Of course
there is plenty added in the typesetting department too. One important new
feature, one that already has been present for a while, is overload protection:
ways to prevent users or third party code to mess up the working of \CONTEXT.

Currently we already enable warnings but we can consider defaulting to error
mode. Why is it needed and how does this impact users? I'll explain some
principles and mention ideas for improvement. We can then discuss and identify
areas that need attention.

And, of course we need to consider how \TEX\ will survive, florish and keep its
niche. Feel free to ask questions, make suggestions and bring in proposals.
